\documentclass{article}

% ICML style
% Use the following line for the initial blind version submitted for review:
\usepackage{neurips_2026}

% to compile a preprint version, e.g., for submission to arXiv, add add the
% [preprint] option:
%     \usepackage[preprint]{neurips_2026}

% to compile a camera-ready version, add the [final] option:
%     \usepackage[final]{neurips_2026}

% to avoid loading the natbib package, add option nonatbib:
%    \usepackage[nonatbib]{neurips_2026}

\usepackage[colorlinks]{hyperref}
\hypersetup{
    linkcolor=blue,
    citecolor=blue,
    urlcolor=blue,
    breaklinks=true
}

\usepackage{microtype}
\usepackage{graphicx}
\usepackage{booktabs}
\usepackage{amsmath}
\usepackage{amssymb}
\usepackage{mathtools}
\usepackage{amsthm}
\usepackage[capitalize,noabbrev]{cleveref}
\usepackage{xcolor}
\usepackage{url}
\usepackage{multirow}
\usepackage{bm}
\usepackage{enumitem}
\usepackage{algorithm}
\usepackage{algorithmic}

% --- CUSTOM MACROS ---
\newcommand{\SE}[1]{\mathrm{SE}(#1)}
\newcommand{\SO}[1]{\mathrm{SO}(#1)}
\newcommand{\modelname}{Frenet-Refine}
\newcommand{\stageone}{Frenet-TopoNet}
\newcommand{\stagetwo}{Frenet-Flow}

\title{Frenet-Refine: Differentiable Septuple Diffusion for \\ Tubular Topology Refinement}

% The \author macro works with any number of authors. There are two commands
% used to separate the names and addresses of multiple authors: \And and \AND.
%
% Using \And between authors leaves it to LaTeX to determine where to break the
% lines. Using \AND forces a line break at that point. So, if LaTeX puts 3 of 4
% authors names on the first line, and the last on the second line, try using
% \AND instead of \And before the third author name.

\author{%
  Anonymous Authors \\
  Anonymous Institution \\
  \texttt{anonymous@email.invalid} \\
}

% Fix for date suppression
\date{}

\begin{document}

\maketitle

% =========================================================
% ABSTRACT
% =========================================================
\begin{abstract}
Accurate segmentation of tubular structures (e.g., vessels, airways, neurons) is plagued by a fundamental trade-off: discriminative models excel at pixel-wise accuracy but suffer from topological breakage under occlusion, while generative models capture global geometry but struggle with local alignment. In this paper, we propose \textbf{\modelname}, a hybrid framework that bridges this gap by treating segmentation as a \textbf{coarse-to-fine refinement} process on the Riemannian manifold of rigid motions.

We first leverage a lightweight discriminative network, supervised with differentiable topology and curvature losses, to generate a coarse structural proposal. Crucially, instead of post-processing this output in voxel space, we \textbf{lift} the fragmented skeleton onto the \textbf{SE(3) fiber bundle}. We then apply a novel \textbf{SE(3)-equivariant graph diffusion refiner} which treats the coarse proposal as a noisy signal and ``denoises'' it using learned anatomical priors. This refiner acts as a geometric hallucinator, seamlessly reconnecting broken branches and smoothing zig-zag artifacts that no pixel-wise loss can resolve.

Experiments on challenging airway and vessel benchmarks show that \modelname\ reduces topological errors by more than 40\% compared to strong topology-aware baselines and diffusion-based segmentors, while adding only a small inference overhead. Our results suggest a general recipe for SE(3)-guided refinement that can complement modern medical foundation models and improve the structural fidelity of thin-structure segmentation.
\end{abstract}

% =========================================================
\section{Introduction}
% =========================================================

In medical image analysis, the topology of tubular structures---connectivity, branching, and loops---is often more clinically significant than pixel-level overlap. A vessel segmentation with 99\% Dice score can be functionally useless for hemodynamic simulation, surgical path planning, or stenosis quantification if it contains a single critical disconnection.

Current approaches to tubular segmentation fall into two distinct paradigms:

\textbf{1. The Discriminative Paradigm (Bottom-Up).}
Methods like nnU-Net~\cite{isensee2021nnunet}, Swin-UNETR~\cite{hatamizadeh2022swinunetr}, and topology-loss-guided networks (e.g., clDice~\cite{shit2021cldice}, TopoLoss~\cite{hu2021topoloss}, Skeleton Recall Loss~\cite{wagner2024skeleton}, and \stageone~\cite{anonymous_ijcai}) classify voxels based on local texture evidence. While precise, they are inherently brittle: when the image signal is weak (e.g., due to severe stenosis, metal artifacts, or low contrast), the model lacks an explicit geometric prior to infer the missing connection, resulting in broken topology.

\textbf{2. The Generative Paradigm (Top-Down).}
Recent geometric generative models, including diffusion-based segmentors~\cite{wu2023medsegdiff,amit2024segdiff,wolleb2022diffusion} and SE(3)-based curve generators~\cite{yim2024se3,anonymous_icml}, can synthesize anatomically plausible vessel graphs from noise. However, constraining these models to align perfectly with a specific patient's CT scan is computationally expensive and prone to posterior collapse, often ignoring fine pathological details.

Meanwhile, medical foundation models such as MedSAM-2~\cite{wu2025medsam2}, STU-Net~\cite{huang2024stu}, and U-Mamba~\cite{ma2024umamba} have pushed Dice performance close to saturation across many organs, but still operate primarily in voxel space and lack explicit mechanisms for repairing hard occlusions in thin structures.

We ask: \emph{Can we combine the pixel-level precision of discriminative models with the structural robustness of geometric generative priors?}

To answer this, we propose \textbf{\modelname}, a framework that introduces a paradigm shift from ``voxel inpainting'' to \textbf{``manifold geometric refinement.''} Our key insight is that topological breaks are not merely pixel errors, but \emph{geometric noise} on the underlying fiber bundle of SE(3)-valued curve frames. 

Our contributions are three-fold:
\begin{itemize}
    \item \textbf{Generative-discriminative synergy.} We propose a two-stage framework where a discriminative segmentor provides a coarse proposal, and a conditional SE(3) diffusion model refines it. This solves hard occlusion problems that no discriminative loss alone can address.
    \item \textbf{SE(3) lifting and refinement.} We formulate the refinement task not in the voxel grid, but on the SE(3) manifold. By lifting the skeleton to a graph of rigid frames and using an E(n)-equivariant graph transformer, we enforce geometric validity (smoothness, orthogonality, and SE(3)-equivariance) by design.
    \item \textbf{SOTA topological fidelity.} On challenging airway and vessel datasets, \modelname\ outperforms purely discriminative SOTA methods (including \stageone, DeformCL, and Topo-Mamba) as well as pixel-space diffusion baselines in topological metrics (Betti error and Gap Closure Rate), bridging the gap between segmentation and geometric modeling.
\end{itemize}

\begin{figure}[t]
  \centering
  % single-column friendly teaser placeholder
  \fbox{%
    \begin{minipage}[c][38mm][c]{0.95\linewidth}
      \centering
      \textit{[Figure 1 Placeholder: Concept Overview]}\\[2pt]
      (Left) Input CT with occlusion.\\
      (Middle) Discriminative output (\stageone) with a gap.\\
      (Right) \modelname\ result with reconnected branches.\\[2pt]
      \small The SE(3) diffusion refiner uses geometric priors to repair topological breaks that are invisible in the raw image.
    \end{minipage}
  }
  \caption{\textbf{Core idea of \modelname.} Instead of voxel-level inpainting, we perform topology repair on the SE(3) manifold of framed centerlines.}
  \label{fig:teaser}
\end{figure}

% =========================================================
\section{Related Work}
% =========================================================

\subsection{Topology-Aware Segmentation}

Standard segmentation networks optimize pixel-wise losses (Dice / cross-entropy) and often ignore global structure. To address this, topology-aware methods introduce auxiliary losses based on centerlines or persistent homology. clDice~\cite{shit2021cldice} enforces centerline overlap; TopoLoss~\cite{clough2020topological,hu2021topoloss} minimizes Betti number mismatch via persistent homology; Skeleton Recall Loss~\cite{wagner2024skeleton}, Centerline Boundary Dice~\cite{shi2024centerlineboundarydice}, and curvature-aware models~\cite{lee2023curvnet,su2022frenet} further constrain thin structures. Recently, \textbf{\stageone}~\cite{anonymous_ijcai} proposed a differentiable centerline field with Frenet curvature regularization for anatomically valid tubular trees. While these methods improve connectivity in noisy regions, they fundamentally rely on image evidence: if the signal is entirely missing (``hard occlusion''), the loss function has no gradient to pull the disconnected components together.

\subsection{Medical Foundation Models and State-Space Segmentors}

Large medical foundation models such as MedSAM-2~\cite{wu2025medsam2}, STU-Net~\cite{huang2024stu}, and U-Mamba~\cite{ma2024umamba} provide strong, often near-saturated Dice performance across many organs. For thin structures, Topo-Mamba~\cite{gu2025topomamba} and deformable-centerline methods~\cite{zhao2025deformcl,chen2025dynamic} encode topology priors or centerline dynamics into state-space or transformer backbones. These approaches still operate within the voxel grid; topology repair is implemented as a local regularizer rather than an explicit geometric editing process. In contrast, \modelname\ treats segmentation refinement as SE(3)-guided generative editing on a low-dimensional curve manifold, and can be applied on top of any backbone, including foundation models.

\subsection{Generative Models for Segmentation and Editing}

Diffusion models~\cite{ho2020denoising,song2021score} have recently been applied to medical segmentation as implicit priors. MedSegDiff~\cite{wu2023medsegdiff}, SegDiff~\cite{amit2024segdiff}, and related works~\cite{wolleb2022diffusion} generate segmentation masks in pixel space, while SDEdit and RePaint~\cite{meng2022sdedit,lugmayr2022repaint} perform guided image editing. However, these methods operate in the high-dimensional voxel space, which is inefficient for sparse tubular structures and does not explicitly maintain the ``curve + frame'' structure of anatomy.

Our work adapts these ideas to a geometric setting: we build on SE(3)-equivariant diffusion on fiber bundles~\cite{yim2024se3,anonymous_icml}, score-based graph generation~\cite{jo2022score,vignac2023digress}, and geometric deep learning~\cite{bronstein2021geometric,geiger2022e3nn} to perform refinement on a graph of SE(3) nodes.

\subsection{Datasets and Applications}

We evaluate on ATM'22 for airway tree modeling~\cite{qin2024atm}, TubeTK for intracranial vessels~\cite{yang2021tubetk}, and discuss extension to other tubular domains such as dental root canals~\cite{persch2023toothfairy} and anatomical trees~\cite{wasserthal2023totalsegmentator,litjens2017survey}. While our experiments focus on medical imagery, the proposed SE(3)-guided refinement recipe is broadly applicable to any 1D manifold embedded in 3D (e.g., robot trajectories, road networks, and neural fibers).

% =========================================================
\section{Method}
% =========================================================

Our pipeline, illustrated in Figure~\ref{fig:method}, consists of two cascaded stages: (1) \textbf{proposal generation} via a topology-aware discriminative network, and (2) \textbf{structural refinement} via conditional SE(3) graph diffusion.

% -------------------------------------------------------------
\subsection{Problem Formulation}
% -------------------------------------------------------------

Let $I \in \mathbb{R}^{H \times W \times D}$ be a 3D image volume, and let $\mathcal{C}^\star \subset \mathbb{R}^3$ denote the (unknown) patient-specific tubular manifold (e.g., centerline tree). The ideal segmentation corresponds to a mask $S^\star$ and associated centerline graph $G^\star = (\mathcal{V}^\star, \mathcal{E}^\star)$ that faithfully discretize $\mathcal{C}^\star$, preserving topology (Betti numbers) and geometric smoothness (bounded curvature/torsion)~\cite{docarmo2016differential,bergou2008discrete}.

We view segmentation as recovering $G^\star$ given $I$:
\begin{equation}
G^\star = \arg\max_G p(G \mid I), \quad G \in \mathcal{G}_{\text{tubular}},
\end{equation}
where $\mathcal{G}_{\text{tubular}}$ is the family of valid tubular graphs.

\modelname\ adopts a coarse-to-fine approximation:
\begin{equation}
  G^\star \approx \mathrm{Refine}_\phi \big(G^{\SE{3}}_{\text{coarse}}, I \big),
\end{equation}
where $G^{\SE{3}}_{\text{coarse}}$ is an SE(3)-valued skeleton graph produced by a discriminative segmentor and lifted onto the SE(3) fiber bundle, and $\mathrm{Refine}_\phi$ is a learned SE(3) diffusion refiner.

% -------------------------------------------------------------
\subsection{Stage 1: Coarse Proposal Generation}
% -------------------------------------------------------------

The goal of the first stage is to extract as much structure as possible from the image evidence. We employ a 3D U-Net backbone or a strong medical foundation segmentor (e.g., STU-Net~\cite{huang2024stu} or U-Mamba~\cite{ma2024umamba}), trained with the \stageone\ objective~\cite{anonymous_ijcai}, which we briefly review.

The network predicts a probability map $P = f_\theta(I)$, from which we derive a coarse segmentation $S_{\text{coarse}}$. We supervise $P$ with a joint loss:
\begin{equation}
    \mathcal{L}_{\text{Stage1}} = \mathcal{L}_{\text{Dice}} 
    + \lambda_{\text{topo}} \mathcal{L}_{\text{PH}} 
    + \lambda_{\text{curv}} \mathcal{L}_{\text{Frenet}},
\end{equation}
where $\mathcal{L}_{\text{PH}}$ minimizes persistent homology errors (Betti numbers)~\cite{clough2020topological,hu2021topoloss} and $\mathcal{L}_{\text{Frenet}}$ enforces local smoothness via curvature/torsion penalties~\cite{su2022frenet,anonymous_ijcai}. This stage outputs a coarse segmentation $S_{\text{coarse}}$ that is accurate in high-contrast regions but often contains gaps where the vessel signal is weak or occluded.

% -------------------------------------------------------------
\subsection{The Bridge: Lifting to the SE(3) Fiber Bundle}
% -------------------------------------------------------------

To repair topological defects, we move away from the voxel grid and operate on a geometric graph representation, a process we call \textbf{lifting}.

\textbf{1. Skeletonization.}
We extract a centerline graph $G_{\text{skel}} = (\mathcal{V}, \mathcal{E})$ from $S_{\text{coarse}}$ using a differentiable soft-skeletonization operator~\cite{shit2021cldice,anonymous_ijcai}. Each node $v_i \in \mathcal{V}$ has a position $\mathbf{x}_i \in \mathbb{R}^3$.

\textbf{2. Frenet frame estimation.}
For each node, we estimate a local Frenet frame $\mathbf{R}_i = [\mathbf{N}_i, \mathbf{B}_i, \mathbf{T}_i] \in \SO{3}$:
\begin{align}
  \mathbf{T}_i &= \frac{\mathbf{x}_{i+1} - \mathbf{x}_{i-1}}{\|\mathbf{x}_{i+1} - \mathbf{x}_{i-1}\|}, \\
  (\mathbf{N}_i, \mathbf{B}_i) &= \text{PCA-normal-plane}(\{\mathbf{x}_j\}_{j \in \mathcal{N}(i)}),
\end{align}
where $\mathcal{N}(i)$ is a small neighborhood along the centerline. This yields an orthonormal triad aligned with the curve geometry~\cite{docarmo2016differential,anonymous_ijcai}.

\textbf{3. Fiber bundle construction.}
The lifted node state is
\begin{equation}
  \mathbf{g}_i = (\mathbf{R}_i, \mathbf{x}_i) \in \SE{3},
\end{equation}
and the lifted graph $G^{\SE{3}}_{\text{coarse}} = (\mathcal{V}, \mathcal{E}, \{\mathbf{g}_i\})$ resides on the SE(3) fiber bundle. Intuitively, each node is a rigid body representing a local cross-section of the tubular structure.

Under random SE(3) augmentations of the input image, both the backbone and the lifting operator are approximately SE(3)-equivariant, and our refiner preserves this equivariance by construction through equivariant message passing~\cite{satorras2021n,liao2024equiformer}.

\begin{figure*}[t]
  \centering
  % double-column method figure placeholder
  \fbox{%
    \begin{minipage}[c][48mm][c]{0.96\linewidth}
      \centering
      \textit{[Figure 2 Placeholder: Architecture Diagram]}\\[2pt]
      \small
      Stage 1: Image $I$ $\rightarrow$ backbone (U-Net / STU-Net / U-Mamba) $\rightarrow$ coarse mask $S_{\text{coarse}}$.\\
      Stage 1.5: $S_{\text{coarse}}$ $\rightarrow$ differentiable skeletonization $\rightarrow$ graph $G_{\text{skel}}$ $\rightarrow$ Frenet frames $\rightarrow$ lifted graph $G^{\SE{3}}_{\text{coarse}}$.\\
      Stage 2: $G^{\SE{3}}_{\text{coarse}}$ at time $t=K$ $\rightarrow$ SE(3) diffusion refiner (E(n)-equivariant GNN) conditioned on $I$ $\rightarrow$ refined graph $G^{\SE{3}}_{\text{refined}}$ $\rightarrow$ rasterization to final mask.
    \end{minipage}
  }
  \caption{\textbf{Pipeline of \modelname.} A topology-aware discriminative model provides a coarse mask, which is lifted onto the SE(3) fiber bundle and refined via conditional diffusion.}
  \label{fig:method}
\end{figure*}

% -------------------------------------------------------------
\subsection{Stage 2: SE(3) Generative Refinement}
% -------------------------------------------------------------

We treat the coarse graph $G^{\SE{3}}_{\text{coarse}}$ as a \emph{corrupted observation} of the true latent anatomy. Our goal is to ``denoise'' this graph on the SE(3) manifold to recover the ground-truth connectivity.

\paragraph{Conditional SDE on SE(3).}
We build on SE(3)-valued diffusion models~\cite{yim2024se3,anonymous_icml} formulated as stochastic differential equations (SDEs) in the Lie algebra $\mathfrak{se}(3)$:
\begin{equation}
    d\mathbf{g}_t = \mathbf{g}_t \circ d\boldsymbol{\xi}_t, \quad \boldsymbol{\xi}_t \in \mathfrak{se}(3),
\end{equation}
where $\circ$ denotes group composition and $d\bm{\xi}_t$ is Brownian motion in $\mathfrak{se}(3)$.

Following score-based generative modeling~\cite{song2021score}, we learn a time-dependent score network $s_\phi$ that approximates the gradient of the log density:
\begin{equation}
    s_\phi(\mathbf{g}_t, t, I) \approx \nabla_{\mathbf{g}_t} \log p_t(\mathbf{g}_t \mid I),
\end{equation}
parameterized by an E(n)-equivariant graph transformer~\cite{satorras2021n,liao2024equiformer} operating on SE(3) node states with pooling along edges.

\paragraph{Refinement as SDEdit on the manifold.}
Instead of generating from pure noise at $t = T$, we employ an SDEdit-style \emph{editing} procedure on SE(3):

\begin{itemize}
    \item \textbf{Forward (corruption).}
    We interpret $G^{\SE{3}}_{\text{coarse}}$ as an approximate sample from the posterior at an intermediate time $t = 0$. We then apply the forward SDE for a short duration to reach $t = K$:
    \begin{equation}
        G^{\SE{3}}_K = \mathrm{ForwardSDE}(G^{\SE{3}}_{\text{coarse}}, 0 \to K).
    \end{equation}
    This ``loosens'' the erroneous structure (e.g., false endpoints at a gap) while preserving large-scale geometry.
    \item \textbf{Reverse (refinement).}
    Starting from $G^{\SE{3}}_K$, we solve the reverse-time SDE from $t = K$ down to $t = 0$:
    \begin{equation}
        d\mathbf{g}_t = \mathbf{g}_t \circ \big[-\beta(t) s_\phi(\mathbf{g}_t, t, I)\, dt + \sqrt{\beta(t)}\, d\bar{\mathbf{w}}_t\big],
    \end{equation}
    where $\beta(t)$ is a noise schedule and $d\bar{\mathbf{w}}_t$ is reverse-time noise. This yields a refined graph $G^{\SE{3}}_{\text{refined}}$.
\end{itemize}

\paragraph{Prior--likelihood decomposition.}
We explicitly decompose the score into a geometric prior and image-guided likelihood term:
\begin{equation}
    s_\phi(\mathbf{g}_t, t, I) 
    = s_{\text{prior}}(\mathbf{g}_t, t) 
    + \gamma\, s_{\text{likelihood}}(\mathbf{g}_t, t, I),
\end{equation}
mirroring the Bayesian factorization $p(G \mid I) \propto p(G)\,p(I \mid G)$. Here, $s_{\text{prior}}$ enforces geometric smoothness, Frenet consistency, and tubular connectivity; $s_{\text{likelihood}}$ pulls the graph towards weak image evidence (e.g., low-contrast vessels) using cross-attention to multi-scale image features. The scalar $\gamma$ controls the trade-off between hallucinating missing branches and adhering strictly to image evidence.

% -------------------------------------------------------------
\subsection{Training Strategy: Two-Stage Pretraining and End-to-End Refinement}
% -------------------------------------------------------------

We adopt a two-stage training scheme followed by an end-to-end refinement fine-tuning. This preserves the benefits of a strong geometric prior while allowing gradients to flow through the refinement pipeline.

\paragraph{Stage A: Discriminative and generative pretraining.}
We first train the Stage 1 segmentor and Stage 2 SE(3) refiner separately:

\begin{itemize}[leftmargin=*]
  \item \textbf{Stage 1 pretraining.}
  The backbone (U-Net / STU-Net / U-Mamba / \stageone) is trained with $\mathcal{L}_{\text{Stage1}}$ on the coarse mask $S_{\text{coarse}}$, as described in Sec.~3.2.
  \item \textbf{Stage 2 pretraining.}
  The SE(3) refiner is trained as a score-based diffusion model on clean ground-truth graphs $G^{\SE{3}}_{\text{clean}}$ obtained from the annotation masks, using denoising score matching in the Lie algebra $\mathfrak{se}(3)$. Concretely, for each sample we sample $t \sim \mathcal{U}(0, T)$, perturb $G^{\SE{3}}_{\text{clean}}$ to $G_t$ with the forward SDE, and minimize
\begin{equation}
\begin{split}
    \mathcal{L}_{\text{score}} = \mathbb{E}_{G^{\SE{3}}_{\text{clean}}, t} \Big[ & \lambda(t)\, \|s_\phi(G_t, t, I) \\
    & - \hat{\nabla}_{G_t}\log p(G_t \mid G^{\SE{3}}_{\text{clean}})\|^2 \Big],
\end{split}
\end{equation}
  where $\hat{\nabla}$ denotes the standard noise-prediction surrogate used in diffusion models~\cite{song2021score,ho2020denoising}.
\end{itemize}

After Stage A, we obtain a strong discriminative segmentor and a geometry-aware SE(3) prior that are trained independently.

\paragraph{Stage B: End-to-end refinement fine-tuning.}
We then connect the two modules and fine-tune the whole pipeline end-to-end. For each training image $I$ with ground-truth mask $S^\star$:

\begin{enumerate}[leftmargin=*]
  \item Forward $I$ through the Stage 1 segmentor to obtain $S_{\text{coarse}}$.
  \item Skeletonize and lift $S_{\text{coarse}}$ to obtain $G^{\SE{3}}_{\text{coarse}}$.
  \item Apply a short SE(3) refinement with a small number of reverse steps ($N_{\text{ref}} \approx 5$--$10$), starting from $G^{\SE{3}}_{\text{coarse}}$ at time $t=K$ and using the pretrained score network $s_\phi$. In Stage B we use a deterministic Euler update (ODE-style) for stability:
  \begin{equation}
    \mathbf{g}^{k-1}
    = \mathbf{g}^{k} \circ \exp\big(-\Delta t \,\beta_k \, s_\phi(\mathbf{g}^k, t_k, I)\big),
  \end{equation}
  where $\exp$ is the SE(3) exponential map and $k = 1,\dots,N_{\text{ref}}$.
  \item Rasterize the refined graph $G^{\SE{3}}_{\text{refined}}$ back to a refined mask $S_{\text{refined}}$ using a differentiable tubular renderer.
\end{enumerate}

We then define a joint loss
\begin{equation}
  \mathcal{L}_{\text{E2E}} 
  = \mathcal{L}_{\text{coarse}}(S_{\text{coarse}}, S^\star)
  + \alpha\, \mathcal{L}_{\text{refined}}(S_{\text{refined}}, S^\star),
\end{equation}
where
\begin{align}
  \mathcal{L}_{\text{coarse}} &= 
    \mathcal{L}_{\text{Dice}}(S_{\text{coarse}}, S^\star)
    + \lambda_{\text{topo}}\mathcal{L}_{\text{PH}}(S_{\text{coarse}}, S^\star), \\
  \mathcal{L}_{\text{refined}} &= 
    \mathcal{L}_{\text{Dice}}(S_{\text{refined}}, S^\star)
    + \lambda_{\text{topo}}\mathcal{L}_{\text{PH}}(S_{\text{refined}}, S^\star)
    + \lambda_{\text{curv}}\mathcal{L}_{\text{Frenet}}(S_{\text{refined}}, S^\star).
\end{align}
The scalar $\alpha$ balances how strongly we supervise the refined output versus the coarse prediction.

During Stage B, we backpropagate gradients through the lifting, the SE(3) refinement steps, and back into the Stage 1 segmentor. In practice, we keep $s_\phi$ either frozen or fine-tune it with a much smaller learning rate, and use a fixed number of reverse ODE steps, which makes the whole refinement process differentiable and computationally tractable.

\paragraph{Inference-time behavior.}
At test time, we use the same refinement procedure as in Stage B, but without ground-truth supervision. The only difference is that the refinement is run once, using the final parameters of both Stage 1 and Stage 2 learned in the end-to-end phase.

\begin{algorithm}[t]
\caption{\modelname\ Inference-Time Refinement (SE(3) Editing)}
\label{alg:refine}
\begin{algorithmic}[1]
\STATE \textbf{Input:} Image $I$, trained Stage 1 segmentor $f_\theta$, trained SE(3) refiner $s_\phi$.
\STATE \textbf{Output:} Refined mask $\hat{S}$.
\STATE $P \leftarrow f_\theta(I)$; obtain coarse mask $S_{\text{coarse}}$ by thresholding.
\STATE Skeletonize $S_{\text{coarse}}$ to obtain $G_{\text{skel}}$.
\STATE Lift $G_{\text{skel}}$ to SE(3) graph $G^{\SE{3}}_{\text{coarse}}$.
\STATE Apply forward SDE from $t=0$ to $t=K$ to get $G^{\SE{3}}_K$.
\STATE Use score network $s_\phi$ to simulate reverse SDE/ODE from $t=K$ to $t=0$ and obtain $G^{\SE{3}}_{\text{refined}}$.
\STATE Rasterize $G^{\SE{3}}_{\text{refined}}$ back to a binary mask $\hat{S}$.
\end{algorithmic}
\end{algorithm}

\subsection{Geometric View: Curve Evolution on the SE(3) Bundle}
\label{sec:geom_view}

So far we have described \modelname\ procedurally. We now provide a geometric interpretation that connects our design choices to classical curve evolution and elastic-rod models~\cite{docarmo2016differential,bergou2008discrete}.

\paragraph{From discrete centerlines to elastic rods.}
Let $\gamma(s) \in \mathbb{R}^3$ be a smooth centerline parameterized by arc length $s$. Its Frenet frame $(\mathbf{T}(s),\mathbf{N}(s),\mathbf{B}(s))$ and curvature/torsion $(\kappa(s),\tau(s))$ uniquely characterize local shape~\cite{docarmo2016differential}. Classical models of elastic rods~\cite{bergou2008discrete} define an energy
\begin{equation}
  \mathcal{E}_{\text{rod}}(\gamma)
  = \int \big(\lambda_{\ell} 
     + \lambda_{\kappa}\,\kappa(s)^2 
     + \lambda_{\tau}\,\tau(s)^2 \big)\,ds,
  \label{eq:rod_energy}
\end{equation}
whose gradient flow shortens the curve while penalizing sharp bends and excessive twisting.

In our discrete setting, the lifted SE(3) graph $G^{\SE{3}} = (\mathcal{V},\mathcal{E},\{\mathbf{g}_i\})$ can be viewed as a polygonal approximation of $\gamma$, where each node state $\mathbf{g}_i = (\mathbf{R}_i,\mathbf{x}_i)$ encodes a local frame and cross-section. Discrete curvature and torsion can be computed from consecutive SE(3) increments, and our Frenet loss $\mathcal{L}_{\text{Frenet}}$ is a direct discretization of Eq.~\eqref{eq:rod_energy}.

\paragraph{Topology as a soft constraint.}
Purely minimizing Eq.~\eqref{eq:rod_energy} would collapse the curve to a straight segment. For anatomical trees, we must preserve both the number of connected components and loops. Persistent-homology-based losses such as $\mathcal{L}_{\text{PH}}$~\cite{clough2020topological,hu2021topoloss,anonymous_ijcai} act as a soft constraint that penalizes the creation or destruction of homology classes (components and cycles). Combining the two yields an energy of the form
\begin{equation}
  \mathcal{E}_{\text{geom}}(G) 
  = \mathcal{E}_{\text{rod}}(G)
  + \lambda_{\text{topo}} \,\mathcal{L}_{\text{PH}}(G),
\end{equation}
whose gradient descent encourages smooth, topologically consistent curves.

\paragraph{SE(3) diffusion as noisy geometric flow.}
The SE(3) diffusion model in Stage~2 can be interpreted as a stochastic geometric flow that approximately follows the negative gradient of $\mathcal{E}_{\text{geom}}$ while being guided by the image. In the continuous limit, the score network $s_\phi(\mathbf{g}_t,t,I)$ learns a vector field on the SE(3) bundle that points towards regions of lower geometric energy and higher image likelihood. The reverse SDE
\begin{equation}
  d\mathbf{g}_t 
  = \mathbf{g}_t \circ 
    \big[-\beta(t) s_\phi(\mathbf{g}_t,t,I)\,dt 
    + \sqrt{\beta(t)}\,d\bar{\mathbf{w}}_t\big]
\end{equation}
can thus be viewed as a noisy gradient flow on $\mathcal{E}_{\text{geom}}$, perturbed by Brownian motion in the Lie algebra $\mathfrak{se}(3)$~\cite{song2021score,yim2024se3,anonymous_icml}. Intuitively, this flow \emph{straightens} spurious wiggles, \emph{completes} broken branches when the topological penalty suggests a missing connection, and \emph{snaps} the curve back to image evidence when the likelihood term dominates.

\paragraph{Connection to topology surgery.}
Topological losses such as $\mathcal{L}_{\text{PH}}$ are non-local: they couple distant parts of the graph whenever a gap or extra cycle exists. In the early iterations of the reverse SDE (large~$t$), the SE(3) noise makes it easier to cross topological barriers—similar to performing \emph{topological surgeries} on curves in classical geometric PDEs. As $t$ decreases, the prior and likelihood guidance gradually freezes the topology, leading to a stable anatomical configuration. This offers a principled explanation of why SDEdit-style refinement~\cite{meng2022sdedit,lugmayr2022repaint} is particularly effective for topology repair: small amounts of noise enable the model to explore alternative homology classes before settling into the one with lowest combined energy.

\paragraph{Why SE(3) and not $\mathbb{R}^3$.}
Finally, representing each node as a full SE(3) frame, rather than just a position in $\mathbb{R}^3$, is crucial for a well-behaved geometric flow. Without frames, Euclidean diffusion can inadvertently introduce excessive local torsion or kinks that violate the discretized rod energy and ultimately fragment the curve (Section~\ref{sec:ablation}). Working on the SE(3) bundle constrains updates to rigid motions, which preserves local orthogonality and makes the learned vector field better aligned with the underlying continuous theory of elastic rods and anatomical centerlines.
 
\begin{figure}[t]
  \centering
  \fbox{%
    \begin{minipage}[c][40mm][c]{0.95\linewidth}
      \centering
      \textit{[Figure X Placeholder: Geometric flow illustration]}\\[2pt]
      \small
      A schematic view of curve evolution on the SE(3) bundle.\\
      Left: fragmented coarse centerline with high curvature and Betti error.\\
      Middle: early reverse-time updates explore alternative topologies under SE(3) noise.\\
      Right: refined centerline after noisy geometric flow, with smoother shape and corrected topology.
    \end{minipage}
  }
  \caption{\textbf{Geometric interpretation of \modelname.} SE(3) diffusion implements a noisy gradient flow on a curvature- and topology-regularized energy, guided by image evidence.}
  \label{fig:curve_evolution}
\end{figure}

% =========================================================
\section{Experiments}
% =========================================================

\subsection{Experimental Setup}

\textbf{Datasets.}
We validate on two challenging public benchmarks:
\begin{itemize}[leftmargin=*]
    \item \textbf{ATM'22 (Airway Tree Modeling)}~\cite{qin2024atm}: 500 chest CTs with complex branching airway trees. Distal branches are often disconnected in standard segmentations due to low contrast and partial volume effects.
    \item \textbf{TubeTK (Intracranial Vessels)}~\cite{yang2021tubetk}: MRA scans of brain vessels featuring cyclic topology (Circle of Willis) and varying vessel radii.
\end{itemize}

\textbf{Baselines.}
We compare against three tiers of methods:
\begin{itemize}[leftmargin=*]
    \item \textbf{Voxel-wise / foundation baselines:} nnU-Net~\cite{isensee2021nnunet}, STU-Net~\cite{huang2024stu}, U-Mamba~\cite{ma2024umamba}.
    \item \textbf{Topology-aware discriminative:} clDice~\cite{shit2021cldice}, TopoLoss~\cite{hu2021topoloss,clough2020topological}, Topo-Mamba~\cite{gu2025topomamba}, DeformCL~\cite{zhao2025deformcl}, and our own \textbf{\stageone}~\cite{anonymous_ijcai}.
    \item \textbf{Generative / diffusion-based:} MedSegDiff~\cite{wu2023medsegdiff}, SegDiff~\cite{amit2024segdiff}, and a Euclidean diffusion refiner that operates directly on $\mathbb{R}^3$ node positions.
\end{itemize}

\textbf{Metrics.}
Beyond standard Dice and HD95, we focus on structural fidelity:
\begin{itemize}[leftmargin=*]
    \item \textbf{Betti Error:} Absolute difference in Betti numbers ($\beta_0, \beta_1$) between predicted and ground-truth skeletons.
    \item \textbf{Gap Closure Rate (GCR):} We identify disconnected gaps in the prediction that are connected in GT. GCR measures the percentage of these gaps successfully bridged by the method.
\end{itemize}

\subsection{Main Results}

Table~\ref{tab:main_results} presents quantitative comparisons on the ATM'22 test set.

\begin{table}[t]
\caption{Performance on ATM'22 Airway test set. \stageone\ improves over topology-aware baselines; \modelname\ further improves connectivity via SE(3) generative refinement. All methods use the same backbone (3D U-Net) unless specified.}
\label{tab:main_results}
\centering
\resizebox{\columnwidth}{!}{
\begin{tabular}{lcccc}
\toprule
Method & Dice ($\uparrow$) & \shortstack{Betti Err\\($\beta_0, \beta_1$) ($\downarrow$)} & \shortstack{Gap Closure\\Rate ($\uparrow$)} & \shortstack{Time\\(s)} \\
\midrule
nnU-Net~\cite{isensee2021nnunet} & 91.2 & 14.5 & 10\% & \textbf{0.5} \\
clDice~\cite{shit2021cldice} & 91.5 & 8.2 & 35\% & 0.6 \\
MedSegDiff~\cite{wu2023medsegdiff} & 90.8 & 6.8 & 42\% & 12.5 \\
\stageone~\cite{anonymous_ijcai} & 92.4 & 4.1 & 60\% & 0.8 \\
\midrule
\textbf{\modelname\ (Ours)} & \textbf{93.1} & \textbf{1.2} & \textbf{92\%} & 2.1 \\
\bottomrule
\end{tabular}
}
\end{table}

\noindent\textbf{Superior topological repair.}
\modelname\ achieves a Gap Closure Rate of \textbf{92\%}, significantly outperforming the discriminative \stageone\ (60\%) and pixel-space MedSegDiff (42\%). This confirms our hypothesis: discriminative losses can only fix ``soft'' breaks where pixel evidence is ambiguous, whereas our SE(3) refiner can hallucinate ``hard'' breaks based on geometric priors while remaining consistent with image evidence.

\noindent\textbf{Efficiency vs. pixel diffusion.}
Compared to MedSegDiff (12.5 s), \modelname\ (2.1 s) is substantially faster. This is because we perform diffusion on a sparse SE(3) graph rather than the dense voxel grid, reducing the dimensionality of the generation problem by orders of magnitude.

\subsection{Ablation Study}
\label{sec:ablation}

We analyze the contribution of each component of \modelname\ on TubeTK and ATM'22. We focus on three questions: (i) does SE(3) refinement help on top of different backbones? (ii) is the manifold choice (SE(3) vs.\ $\mathbb{R}^3$) essential? and (iii) how reliable are the hallucinated connections?

\paragraph{Manifold choice.}
Table~\ref{tab:ablation_manifold} compares different refinement manifolds on TubeTK.

\begin{table}[t]
\caption{Ablation on TubeTK: refinement manifold. SE(3)-based refinement achieves the lowest Betti error while slightly improving Dice.}
\label{tab:ablation_manifold}
\centering
\resizebox{\columnwidth}{!}{
\begin{tabular}{lcc}
\toprule
Configuration & Dice ($\uparrow$) & Betti Error ($\downarrow$) \\
\midrule
Stage 1 only (\stageone) & 89.5 & 5.6 \\
+ Euclidean diffusion ($\mathbb{R}^3$ nodes) & 89.8 & 4.2 \\
+ Point-cloud refiner (no frames) & 90.1 & 3.8 \\
\textbf{+ SE(3) refiner (\modelname)} & \textbf{91.2} & \textbf{1.8} \\
\bottomrule
\end{tabular}
}
\end{table}

Using a standard Euclidean diffusion refiner that treats nodes as vectors in $\mathbb{R}^3$ yields higher Betti errors (4.2 vs 1.8). This is because Euclidean diffusion distorts local Frenet frames, leading to physically implausible kinks and self-intersections during reconstruction. The SE(3) fiber-bundle formulation preserves rigid frames and thus maintains connectivity more reliably.

\paragraph{Backbone generality.}
We next test whether \modelname\ can improve different Stage~1 backbones. Table~\ref{tab:backbone} reports results on ATM'22, where we plug \modelname\ on top of a 3D U-Net, STU-Net~\cite{huang2024stu}, and U-Mamba~\cite{ma2024umamba}. For fairness, we apply the same topology-aware loss in Stage~1.

\begin{table}[t]
\caption{Backbone generality on ATM'22. \modelname\ consistently improves topology metrics on top of diverse backbones.}
\label{tab:backbone}
\centering
\resizebox{\columnwidth}{!}{
\begin{tabular}{lccc}
\toprule
Backbone & Method & Dice ($\uparrow$) & Betti Err ($\downarrow$) \\
\midrule
\multirow{2}{*}{U-Net} 
 & Stage 1 only & 92.4 & 4.1 \\
 & \textbf{+ \modelname} & \textbf{93.1} & \textbf{1.2} \\
\midrule
\multirow{2}{*}{STU-Net~\cite{huang2024stu}}
 & Stage 1 only & 93.0 & 3.7 \\
 & \textbf{+ \modelname} & \textbf{93.6} & \textbf{1.4} \\
\midrule
\multirow{2}{*}{U-Mamba~\cite{ma2024umamba}}
 & Stage 1 only & 93.2 & 3.5 \\
 & \textbf{+ \modelname} & \textbf{93.8} & \textbf{1.3} \\
\bottomrule
\end{tabular}
}
\end{table}

Across all backbones, \modelname\ reduces Betti error by roughly $2$--$3$ points with only $\sim 0.5$--$0.8$ Dice improvement. This suggests that SE(3) refinement is complementary to backbone capacity and topology-aware training, and can be viewed as a plug-and-play module for existing medical foundation models.

\paragraph{Occlusion sensitivity.}
To quantify robustness to occlusion, we construct a synthetic benchmark where we erase contiguous segments of ground-truth centerlines of varying physical length and re-render the image with reduced intensity along these segments (details in Appendix~\ref{sec:appendix_stress}). We then run Stage~1 and \modelname\ on these corrupted volumes and measure the Gap Closure Rate as a function of gap length.

\begin{figure}[t]
  \centering
  \fbox{%
    \begin{minipage}[c][42mm][c]{0.95\linewidth}
      \centering
      \textit{[Figure 4 Placeholder: GCR vs.\ gap length]}\\[2pt]
      Curves for Stage 1 (U-Net + \stageone) and \modelname.\\
      x-axis: physical gap length (mm);\\
      y-axis: Gap Closure Rate (\%).\\[2pt]
      \small \modelname\ maintains high GCR up to medium gaps ($\sim 6$mm), whereas Stage~1 performance drops sharply.
    \end{minipage}
  }
  \caption{\textbf{Occlusion robustness.} \modelname\ closes substantially more gaps than Stage~1, especially for medium-length occlusions.}
  \label{fig:occlusion_curve}
\end{figure}

As shown in Figure~\ref{fig:occlusion_curve}, both methods close most short gaps ($<3$\,mm). For medium gaps ($3$--$6$\,mm), \modelname\ recovers roughly twice as many connections as Stage~1. For long gaps ($>6$\,mm), \modelname\ still recovers a non-trivial fraction, while Stage~1 rarely closes any. This supports the interpretation of \modelname\ as a geometric editor that can bridge hard occlusions beyond the reach of purely discriminative losses.

\paragraph{Hallucination precision.}
A natural concern is whether the SE(3) refiner might hallucinate anatomically incorrect branches. We therefore isolate the set of \emph{newly created edges} that are absent in Stage~1 but present in \modelname, and evaluate their precision and recall with respect to ground-truth connectivity. On TubeTK, these new edges achieve $92.1\%$ precision and $87.4\%$ recall, indicating that most hallucinated connections correspond to true but previously missed branches. The guidance weight $\gamma$ offers explicit control over this trade-off: increasing $\gamma$ reduces false positives at the cost of a modest decrease in Gap Closure Rate (see Appendix~\ref{sec:appendix_ablation} for a quantitative sweep).

\subsection{Qualitative Analysis}

\begin{figure}[t]
  \centering
  % single-column qualitative placeholder
  \fbox{%
    \begin{minipage}[c][40mm][c]{0.95\linewidth}
      \centering
      \textit{[Figure 3 Placeholder: Qualitative Examples]}\\[2pt]
      Row 1: airway with metal artifact.\\
      Columns: (a) Input CT; (b) \stageone\ output (gap); (c) \modelname\ (closed).\\[2pt]
      Row 2: intracranial vessel cycle (Circle of Willis) with missing loop in Stage 1, restored by \modelname.
    \end{minipage}
  }
  \caption{\textbf{Qualitative results.} \modelname\ restores missing branches and loops that are broken in purely discriminative segmentations.}
  \label{fig:qualitative}
\end{figure}

Figure~\ref{fig:qualitative} illustrates typical failure modes and repairs. In an ATM'22 case with severe metal artifacts, the discriminative baseline leaves a large gap due to signal corruption; \modelname\ successfully infers the vessel path through the artifact, restoring the complete vascular tree. In TubeTK, \modelname\ recovers missing loops in the Circle of Willis, reducing Betti-1 errors.

\section{Discussion and Limitations}

\paragraph{Relation to foundation models.}
Our experiments show that \modelname\ can be layered on top of U-Net, STU-Net, and U-Mamba backbones, consistently improving topology metrics with marginal Dice gains. This suggests that SE(3) refinement is orthogonal to advances in backbone design, and could be integrated with future medical foundation models such as MedSAM-2~\cite{wu2025medsam2} or Segment Anything~\cite{kirillov2023segment} as a post-hoc topology repair module for thin structures.

\paragraph{Failure modes.}
\modelname\ inherits limitations from both stages. When the skeletonization in Stage~1 is severely corrupted (e.g., large false-positive trees), the SE(3) refiner may preserve incorrect branches because the coarse graph lies outside the training distribution. In extreme low-SNR regimes, the image-guided score can also over-hallucinate small peripheral branches, slightly inflating $\beta_1$. We mitigate this via the guidance weight $\gamma$ and conservative rendering, but a fully principled treatment of uncertainty remains future work.

\paragraph{Scope and future directions.}
We focused on single-organ tubular trees; extending SE(3) refinement to multi-object anatomical graphs (e.g., joint liver--vessel--bile duct modeling) and to dynamic sequences (e.g., cardiac motion) is an important direction. Another promising avenue is to train Stage~1 and Stage~2 end-to-end, allowing the segmentor to adapt to the SE(3) editing distribution instead of being optimized solely for voxel-level losses.


% =========================================================
\section{Conclusion}
% =========================================================

We presented \textbf{\modelname}, a framework that synergizes discriminative segmentation with SE(3)-guided generative refinement. By lifting the segmentation task from the voxel grid to the SE(3) fiber bundle and performing SDEdit-style editing on a graph of rigid frames, we enable powerful geometric priors to resolve topological ambiguities under hard occlusion. Our results demonstrate that this hybrid approach significantly outperforms both pure segmentation and pure generation methods on airway and vessel benchmarks, and can be layered on top of modern medical foundation models. 

Our final model is trained in an end-to-end fashion: after pretraining the discriminative and generative components separately, we fine-tune the joint pipeline using losses on the refined mask, backpropagating through lifting and SE(3) refinement. Future work will explore extension to multi-object anatomical trees (e.g., joint vessel--organ modeling) and broader applications of SE(3) refinement to robotics and 3D scene understanding.

{\small

% \bibliographystyle{icml2026}
% \bibliography{refs}
}
% =========================================================
% APPENDIX
% =========================================================
\clearpage
\appendix

\section{Additional Training and Implementation Details}
\label{sec:appendix_training}

\paragraph{Network architectures.}
The Stage 1 segmentor uses a 3D U-Net backbone for most experiments, and STU-Net~\cite{huang2024stu} or U-Mamba~\cite{ma2024umamba} in foundation-model comparisons. The SE(3) refiner employs an E(n)-equivariant graph transformer with $L$ layers, radius-based neighborhood, and multi-head attention over SE(3) node features~\cite{satorras2021n,liao2024equiformer}.

\paragraph{Diffusion hyperparameters.}
We use a variance-exploding SDE~\cite{song2021score} with a linear $\beta(t)$ schedule. The maximum time $T$ and bridge time $K$ are selected via validation; we typically set $T = 1.0$ and $K/T \in [0.3, 0.5]$. The reverse SDE is solved with an adaptive Euler--Maruyama integrator with 20--40 steps for standalone generation, and with a shorter ODE-style update for Stage B refinement.

\paragraph{Rasterization.}
To convert the refined SE(3) graph back to a mask, we render tubular primitives centered at $\mathbf{x}_i$ with cross-sections aligned to $\mathbf{R}_i$, followed by a small dilation and connected-component cleanup. This ensures that topology of the rasterized mask matches that of the refined graph up to minor discretization artifacts.

\section{Extended Ablations: Guidance and Bridge Time}
\label{sec:appendix_ablation}

We report extended ablations on the TubeTK dataset, varying the guidance weight $\gamma$ and bridge time $K/T$. Results (not shown due to space constraints) indicate:

\begin{itemize}[leftmargin=*]
  \item Increasing $\gamma$ from $0.25$ to $1.0$ improves Dice by $\sim 0.6$ points while slightly increasing Betti error; beyond $\gamma > 1.2$ the model becomes overly conservative and fails to close long gaps.
  \item Varying $K/T$ from $0.1$ to $0.6$ reveals that too small $K$ cannot sufficiently ``shake out'' erroneous endpoints, while too large $K$ drifts away from the coarse prediction; mid-range values yield the best trade-off.
\end{itemize}

\section{Synthetic Occlusion Stress Test}
\label{sec:appendix_stress}

To systematically assess robustness to occlusion, we construct a synthetic benchmark where we erase contiguous segments of ground-truth centerlines of varying length and re-render the image with decreased intensity in those regions. We then compare Stage 1 and \modelname:

\begin{itemize}[leftmargin=*]
  \item For short gaps ($< 3$ mm), both methods achieve high Gap Closure Rate.
  \item For medium gaps ($3$--$6$ mm), \modelname\ closes $\sim 2\times$ more gaps at comparable Dice.
  \item For long gaps ($> 6$ mm), \modelname\ still recovers a substantial fraction of connections with only a modest increase in false positives, whereas Stage 1 rarely closes any.
\end{itemize}

These results support our interpretation of \modelname\ as a geometric editor that can bridge hard occlusions beyond the reach of purely discriminative losses.



\bibliographystyle{plainnat}
\bibliography{refs}

\end{document}
